% !PRINTED 610
\weberpage{610}
%<<CONT

1. \gesperrt{Ist $\Omega$ ein quadratischer Körper mit der Grundzahl $\varDelta$, und $\varphi$ ein ganzes Functional in diesem Körper, so lässt sich eine zur Discriminante $\varDelta$ gehörige quadratische Irrationalzahl}

\weq{17}{\omega = \frac{b + \sqrt{\varDelta}}{2c}}

\gesperrt{und eine positive ganze rationale Zahl $e$ so bestimmen, dass}

\weq{18}{ec(t_1 + \omega t_2)}

\gesperrt{eine Basisform von $\varphi$ wird. Umgekehrt ist, wenn $\omega$ eine beliebige zur Discriminante $\varDelta$ gehörige quadratische Irrationalzahl ist, die Form (18) immer eine Basisform eines Functionals $\varphi$.}

\begin{center}
§. 167.

\gesperrt{Aequivalente Functionale und äquivalente Zahlen.}
\end{center}

Wenn nun $\omega, \omega_1$ irgend zwei zu der Discriminante $\varDelta$ gehörige quadratische Irrationalzahlen sind, so erhalten wir zwei Basisformen

\weq{1}{\begin{aligned} \varphi &= ec(t_1+\omega t_2)\\ \varphi_1 &= e_1c_1(t_1+\omega_1 t_2). \end{aligned}}

Wenn nun die beiden Zahlen $\omega, \omega_1$ in dem Sinne äquivalent sind, wie wir diesen Begriff im \S.\,120 des ersten Bandes festgesetzt haben, so ist

\weq{2}{\omega_1 = \frac{p\omega+q}{r\omega+s},}

worin $p,q,r,s$ ganze rationale Zahlen sind, die der Bedingung

\weq{3}{ps - qr = \pm 1}

genügen. Machen wir dann die Substitution (2) in $\varphi_1$, so folgt

\[ (r\omega+s)\varphi_1 = e_1c_1[st_1+qt_2+\omega(rt_1+pt_2)], \]

oder wenn wir

\weq{4}{u_1 = st_1+qt_2,\quad u_2 = rt_1+pt_2}

setzen,

\[ (r\omega+s)\varphi_1 = e_1c_1(u_1+\omega u_2). \]

Nun ist $ec(u_1+\omega u_2)$ mit $\varphi$ associirt, weil (4) eine ganz-
%CONT>>
